\renewcommand*{\authorOfNextLines}{Helmut Quinz}
\TeX ~ist ein von Donald E. Knuth ab 1977 entwickeltes und 1986 fertiggestelltes Textsatzsystem.
Hauptfunktion ist das grafisch ansprechende komponieren von Text und Grafiken.
Das \TeX -Grundpaket verfügt über eine Makrosprache um den Text-Satz beeinflussen zu können.
Die Entwicklung des \TeX -Grundpaketes gilt, abgesehen von der Behebung unerwünschten Verhaltens als abgeschlossen.
\TeX wird nicht zuletzt auf Grund der guten Unterstützung bei der Darstellung mathematischer Zusammenhänge im technisch-universitären Bereich verwendet.
Das Textsatzsystem wird von einer akttiven Gemeinschaft getragen die für dieses System eine Unzahl an Makrosammlungen bereitstellen und weiterentwickeln.
Diese Makrosamlungen sind wie das \TeX -Grundpaket frei verfügbar. 
Heute wird \TeX ~üblicherweise eingebunden in leistungsfähige Erweiterungspakete verwendet. 
Bekanntestes, wenn auch nicht alleinstehendes Beispiel dafür ist \LaTeX.
\TeX ~ist rechnerunabhängig verwendbar und sollte auf allen Betriebssystemen vergleichbare Ergebnisse liefern.\\

\TeX ~wird heute von verschiedenen Distributoren zur Verfügung gestellt.
Dazu gehören
\begin{itemize}
	\item MiKTeX,
	\item TeX Live oder
	\item teTex (wird seit Mai 2006 nicht mehr weiterentwickelt) 
\end{itemize}

Die Pakete stellen jeweils unterschiedliche Compiler-Sätze bereit, welche sich aus der historischen Entwicklung des Textsatzsystems erklären lassen.
so wurde LaTeX, welches in der Grundeinstellung \texttt{.DVI} Dateien erstellt zu PdfLaTeX weiterentwickelt.
Letzteres wurde unter Berücksichtigung aktueller Entwicklungen (utf8 als Standart-Zeichensatz) zu LuaLaTex und XeLatTex erweitert.